\section[Conclusão]{Conclusão}

A tragédia que assolou o Rio Grande do Sul impactou diretamente o fluxo de trabalho, engajamento e andamento do projeto. Além disso, a falta de engajamento e de críticas construtivas por parte das \textit{stakeholders} foi desmotivadora, especialmente quando apenas uma delas comparecia às apresentações e não fornecia nenhum tipo de \textit{feedback}.

No entanto, a experiência de desenvolver um projeto com entregas iterativas e incrementais, com o compromisso de entregar um sistema funcional ao final, foi muito valiosa. A cada sprint, enfrentamos desafios como a integração das funcionalidades e a refatoração necessária para resolver problemas técnicos, o que nos permitiu aprimorar nossas habilidades.

Enfrentamos alguns problemas de comunicação dentro da equipe, resultando em atrasos na entrega de algumas funcionalidades. O principal obstáculo foi a necessidade de refatorar a integração de imagens devido a mudanças no \textit{backend} e dificuldades com a \ac{aws}. Apesar desses desafios, conseguimos entregar um produto funcional. Sinto que poderíamos ter alcançado um resultado ainda melhor se não fossem os imprevistos e a falta de engajamento das \textit{stakeholders}.

No fim deste projeto, me sinto um pouco frustrado, pois acredito que poderíamos ter entregue um projeto muito melhor caso não houvesse ocorrido esses pontos críticos. No entanto, este semestre me proporcionou um crescimento pessoal único. Saber lidar com esses tipos de obstáculos durante o desenvolvimento de software é um aprendizado incrível para que eu possa melhorar minhas habilidades técnicas e de resolução de problemas.

De forma geral, foi muito satisfatório participar do desenvolvimento da aplicação Decola e espero que o trabalho desenvolvido na \ac{ages} seja continuado pelas clientes.